\section{General scalar entropy-SGD
exit}\label{scalar_exit-general-scalar-entropy-sgd-exit}

Part of \href{DRAFT.md}{the manuscript}. E-numbering and section numbers
are continuous across \href{scalar_exit.md}{scalar exit},
\href{calibration.md}{calibration}, and \href{risk_boundaries.md}{risk
boundaries}. Full statements and proofs are retained. Constants are not
numerically calibrated; integration is not a completed independent proof
audit. Use \(B_U=w/4\) for the update bound and \(M\) for calibration
sample size; the original Gaussian sections use \(M=w/4\) locally.

\subsection{1. The process being
predicted}\label{scalar_exit-the-process-being-predicted}

Let \(Q\) be a fixed probability distribution on scalar scores
\(Z\in\mathbb R\). It may be continuous, discrete, or a fixed weighted
empirical distribution. Fix \(w>0\), define \[
s(u)=(1+e^{-u})^{-1},\qquad a(u)=s(u)(1-s(u)),
\] and use the actual threshold entropy velocity \[
U(x,z)=-w^2(z-x)a(w(z-x)),\qquad
v(x)=\mathbb E_Q U(x,Z).
\tag{E.1}
\] The population map and fresh-batch update are \[
G(x)=x+\eta v(x),\qquad
x_{t+1}=x_t+\frac{\eta}{n}\sum_{i=1}^n U(x_t,Z_{t,i}),
\tag{E.2}
\] where the \(Z_{t,i}\) are iid from \(Q\), and the initial state is
deterministic unless stated otherwise.

An independent calibration sample \[
\widetilde Z_1,\ldots,\widetilde Z_M\overset{\mathrm{iid}}{\sim}Q
\tag{E.3}
\] will be used only to construct a predictor. It does not alter
\(x_0\), \(Q\), the adaptation samples, or the recursion (E.2).
Conditioning on a fixed feature pool, its source projection, and any
externally fixed design parameters is permitted. All probabilities in
that case are conditional on those objects.

For intervals \(I_n=[a_n,c_n]\), define \[
\tau_n=\inf\{t\ge0:x_t\notin I_n\},\qquad
J_n=\begin{cases}+1,&x_{\tau_n}>c_n,\\-1,&x_{\tau_n}<a_n.\end{cases}
\tag{E.4}
\] The first part of the analysis concerns this interval-exit event. It
becomes a classification-risk exit theorem when the relevant risk-safe
set is exactly \(I_n\), or when the event discrepancy is separately
controlled as in \href{risk_boundaries.md}{Section 7}.

\subsubsection{Lemma E.1. Distribution-free regularity of the entropy
update}\label{scalar_exit-lemma-e.1.-distribution-free-regularity-of-the-entropy-update}

For every \(x,z\), \[
|U(x,z)|\le B_U:=w/4,\qquad
|\partial_xU(x,z)|\le L_U:=w^2/4.
\tag{E.5}
\] Every higher fixed-order derivative in \(x\) is also uniformly
bounded. In particular, \(|\partial_x^2 U|\le 3w^3/4\). Consequently
\(v\) is smooth for every probability distribution \(Q\), without moment
assumptions on \(Z\).

If \(0<\eta<1/L_U\), both \(G\) and every empirical population map
\(x+\eta M^{-1}\sum_i U(x,\widetilde Z_i)\) are increasing smooth
bijections of the line.

\textbf{Proof.} If \(h\) is the binary entropy of \(s\), then
\(h'(u)=-ua(u)\) and \[
h''(u)=a(u)\{u\tanh(u/2)-1\}.
\] To verify the constants, write \(a(u)=[2(1+\cosh u)]^{-1}\). The
inequality \(1+\cosh u\ge2+u^2/2\ge2|u|\) gives \(|u|a(u)\le1/4\). The
two nonnegative terms \(ua(u)\tanh(u/2)\) and \(a(u)\) each lie in
\([0,1/4]\), proving \(|h''|\le1/4\). Differentiation gives \[
h'''(u)=2a(u)\tanh(u/2)+ua(u)\{6a(u)-1\},
\] whose absolute value is at most \(1/2+1/4=3/4\). These bounds give
(E.5) and the second-derivative bound after scaling. Further derivatives
are sums of polynomial factors times derivatives of the logistic
function, and decay at both infinities; hence they are bounded.
Dominated differentiation applies under any \(Q\). Finally, \[
G'(x)\ge1-\eta L_U>0,\qquad |G(x)-x|\le\eta B_U.
\] Strict monotonicity and the limits at infinity give bijectivity. The
empirical map satisfies the same inequalities. \(\square\)

\subsection{2. General scalar exit
geometry}\label{scalar_exit-general-scalar-exit-geometry}

Fix a compact interval \(\mathcal K=[\ell,r]\) and \(q\in(\ell,r)\). We
impose the following conditions.

\textbf{G1: one repeller on the required domain.} \[
v(q)=0,\qquad \lambda:=v'(q)>0,\qquad
(x-q)v(x)>0\quad(x\in\mathcal K,\ x\ne q).
\tag{E.6}
\] No uniqueness of the entropy peak outside \(\mathcal K\) is required.

\textbf{G2: nondegenerate noise.} \[
\sigma_g^2:=\operatorname{Var}_Q U(q,Z)>0.
\tag{E.7}
\]

\textbf{G3: fixed step size and interior exit boundaries.} Fix
\(0<\eta<1/L_U\), and set \[
A=1+\eta\lambda>1,\qquad
s^2=\frac{\eta^2\sigma_g^2}{A^2-1}.
\tag{E.8}
\] For fixed \(d_0,\gamma>0\), all sufficiently large \(n\) satisfy \[
a_n\le q-d_0,\qquad c_n\ge q+d_0,
\] \[
[a_n-\eta B_U,c_n+\eta B_U]
\subset[\ell+\gamma,r-\gamma].
\tag{E.9}
\] Thus both the deterministic proposal and the actual one-step
overshoot from \(I_n\) remain in the domain where the geometry is
controlled.

\textbf{G4: critical initialization.} \[
|x_{0,n}-q|\le K_0/\sqrt n
\tag{E.10}
\] for a fixed finite \(K_0\). In particular \(x_{0,n}\in I_n\) for
large \(n\).

Constants below may depend on \(Q,w,\eta,\mathcal K,d_0,\gamma,K_0\).
They do not depend on \(n,t,M\). Uniformity over a family of score
distributions additionally requires uniform versions of these geometric
and noise bounds; it is not automatic when a curvature or an outward
drift bottleneck approaches zero.

\subsubsection{Lemma E.2. Normalized conjugacy on the exit
domain}\label{scalar_exit-lemma-e.2.-normalized-conjugacy-on-the-exit-domain}

Under G1 and G3, there is an increasing \(C^2\) function \(H\) on
\(\mathcal K\), with bounded derivatives and a strictly positive lower
bound on \(H'\), such that \[
H(q)=0,\qquad H'(q)=1,\qquad H(G(x))=A H(x)
\tag{E.11}
\] whenever \(x,G(x)\in\mathcal K\). It is the unique normalized \(C^1\)
conjugacy on this domain.

\textbf{Proof.} Since \(G\) is increasing and \(G(q)=q\), its inverse
sends each point of \(\mathcal K\) toward \(q\). For example, if
\(x>q\), then \[
q<G^{-1}(x)<x,
\] using \(G(x)>x\) and monotonicity. The inverse orbit stays in
\(\mathcal K\) and converges to \(q\); any other limit would be another
zero of \(v\).

Near zero define the centered inverse \[
g(y)=G^{-1}(q+y)-q=A^{-1}y+O(y^2).
\] Choose a sufficiently small neighborhood and \(b<1\) such that \[
\sup |g'|\le b,\qquad Ab^2<1.
\] Writing \(R(y)=g(y)-A^{-1}y\), the approximants \[
h_j(y)=A^j g^j(y)
\] satisfy \[
h_{j+1}(y)-h_j(y)=A^{j+1}R(g^j(y)).
\tag{E.12}
\] Here \(g^j\) denotes composition. Iterated differentiation gives \[
\|g^j\|_{C^2}\le Cb^j
\] on this local interval. Since \(R(0)=R'(0)=0\), the \(C^2\) norm of
the right-hand side of (E.12) is at most \(C(Ab^2)^j\). The sequence
converges in \(C^2\) to \(h\), with \(h(0)=0,h'(0)=1\), and \(h'>0\)
after decreasing the local interval if necessary.

A fixed finite number \(N\) of inverse steps brings all of
\(\mathcal K\) into that local interval. This follows first for the two
endpoints and then for every interior point by monotonicity. Set \[
H(x)=A^N h(G^{-N}(x)-q).
\tag{E.13}
\] The local functional equation makes this definition independent of a
sufficiently large \(N\). It gives (E.11), regularity, and positivity of
\(H'\). Compactness supplies the derivative bounds.

For uniqueness, any other normalized \(C^1\) conjugacy \(\widetilde H\)
satisfies \[
\widetilde H(x)=A^j\widetilde H(G^{-j}(x)).
\] Since \(\widetilde H(q+y)=y+o(y)\) and \(A^j(G^{-j}(x)-q)\to H(x)\),
the right-hand side tends to \(H(x)\). \(\square\)

\subsubsection{Theorem E.3. Joint exit approximation for general scalar
score
laws}\label{scalar_exit-theorem-e.3.-joint-exit-approximation-for-general-scalar-score-laws}

Under G1--G4, \(\tau_n<\infty\) almost surely. Define \[
\mu_n=\sqrt n H(x_{0,n}),\qquad
B_n^+=H(c_n),\qquad B_n^-=-H(a_n),
\tag{E.14}
\] and \[
\Psi_n^+(t)=
\Phi\!\left(\frac{\mu_n-\sqrt n A^{-t}B_n^+}{s}\right),
\qquad
\Psi_n^-(t)=
\Phi\!\left(\frac{-\mu_n-\sqrt n A^{-t}B_n^-}{s}\right).
\tag{E.15}
\] The \(B_n^\pm\) are uniformly bounded above and away from zero. For
all integers \(t\ge0\), \[
\left|\Pr(\tau_n\le t,J_n=\pm1)-\Psi_n^\pm(t)\right|
\le C\left[n^{-1/2}+
\left(\frac{1+\log n}{n}+A^{-2t}\right)^{1/3}\right].
\tag{E.16}
\] Moreover, with \[
e_n=n^{-1/2}+\left(\frac{1+\log n}{n}\right)^{1/3},
\tag{E.17}
\] the stronger uniform-in-time statement holds: \[
\sup_{t\in\mathbb Z_{\ge0}}
\max_{\pm}\left|\Pr(\tau_n\le t,J_n=\pm1)-\Psi_n^\pm(t)\right|
\le Ce_n.
\tag{E.18}
\] In particular, \[
\left|\Pr(J_n=+1)-\Phi(\mu_n/s)\right|\le Ce_n.
\tag{E.19}
\] The same conclusions hold for open or half-open intervals, with their
declared endpoint membership used consistently in the exit event and in
\(J_n\).

\textbf{Proof.}

\textbf{Step 1: finite exit without Gaussian tails.} Since
\(\mathbb E U(q,Z)=0\) and its variance is positive, both \(U(q,Z)>0\)
and \(U(q,Z)<0\) have positive probability. The sign of \(U(q,z)\) is
the sign of \(q-z\). There therefore exist bounded score intervals
strictly to the left and right of \(q\), with positive masses
\(p_L,p_R\). For states \(x\in[q,c_n]\), samples from the left interval
produce a uniformly positive velocity; for states \(x\in[a_n,q]\),
samples from the right interval produce a uniformly negative velocity.
Uniformity follows from continuity and compactness, and from the
positive separation of the selected score intervals from \(q\).

Choose a fixed block length \(N_0\) large enough that either such
sequence of updates exits the relevant side. At the start of any block
before exit, select the left-score event if the state is at least \(q\),
and the right-score event otherwise. The conditional probability of the
selected event is at least \[
p_n:=\min(p_L,p_R)^{nN_0}>0.
\] Thus \(\Pr(\tau_n>kN_0)\le(1-p_n)^k\). This argument applies to
discrete score laws as well.

\textbf{Step 2: stopped conjugacy coordinates.} Write \[
\xi_j(x)=\frac1n\sum_{i=1}^n U(x,Z_{j,i})-v(x).
\] Conditional on the history before batch \(j\), for predictable \(x\),
\[
\mathbb E\xi_j(x)^2\le B_U^2/n,\qquad
\mathbb E\xi_j(x)^4\le C/n^2,
\] \[
\mathbb E(\xi_j(x)-\xi_j(q))^2\le L_U^2|x-q|^2/n,
\qquad
\mathbb E e^{u\xi_j(x)}\le e^{u^2B_U^2/(2n)}.
\tag{E.20}
\] Taylor expansion on the overshoot domain gives, before exit, \[
H(x_{j+1})=AH(x_j)+
\eta H'(G(x_j))\xi_j(x_j)+R_j,
\qquad |R_j|\le C\xi_j(x_j)^2.
\tag{E.21}
\] Define \[
Y_{n,t}=\sqrt n A^{-(t\wedge\tau_n)}
H(x_{t\wedge\tau_n}).
\] Summation of (E.21) expresses \(Y_{n,t}\) as \(\mu_n\), a stopped
martingale \(\mathcal M_{n,t}\), and a remainder \(\mathcal E_{n,t}\).
Orthogonality and (E.20) yield \[
\sup_t\mathbb E|\mathcal M_{n,t}|^2\le C,\qquad
\sup_t\|\mathcal E_{n,t}\|_2\le C/\sqrt n,
\qquad \sup_t\mathbb E|Y_{n,t}|^2\le C.
\tag{E.22}
\] For the remainder, use Minkowski's inequality and
\(\|\mathbf1_{\{j<\tau_n\}}\xi_j(x_j)^2\|_2\le C/n\) in the
geometrically weighted sum. Fatou's lemma also gives the corresponding
estimate at exit. Since \(|H(x_{\tau_n})|\ge b_0>0\), \[
\mathbb E A^{-2\tau_n}\le C/n.
\tag{E.23}
\]

\textbf{Step 3: freeze the noise at the actual root.} On the same
infinite independent sample array, define \[
W_n=\mu_n+\eta\sum_{j\ge0}A^{-j-1}\sqrt n\,\xi_j(q).
\tag{E.24}
\] Before exit, \[
D_j=H'(G(x_j))\xi_j(x_j)-\xi_j(q)
\] has conditional mean zero and conditional variance at most
\(C|x_j-q|^2/n\). The lower bound on \(H'\), (E.22), and compactness
imply \[
\mathbb E[\mathbf1_{\{j<\tau_n\}}|x_j-q|^2]
\le C\min(A^{2j}/n,1).
\tag{E.25}
\] Consequently, \[
\mathbb E\left|
\sqrt n\eta\sum_{j<t\wedge\tau_n}A^{-j-1}D_j
\right|^2
\le C\sum_{j\ge0}A^{-2j}\min(A^{2j}/n,1)
\le C(1+\log n)/n.
\tag{E.26}
\] The frozen-noise tail after \(t\wedge\tau_n\) has conditional
variance \(s^2A^{-2(t\wedge\tau_n)}\). Equation (E.23) bounds its
expected contribution. Together with the remainder, \[
\mathbb E|Y_{n,t}-W_n|^2
\le C\left[\frac{1+\log n}{n}+A^{-2t}\right],
\] \[
\mathbb E|Y_{n,\tau_n}-W_n|^2
\le C(1+\log n)/n.
\tag{E.27}
\] This comparison controls the actual state-dependent batch updates.

\textbf{Step 4: Gaussian seed and event identities.} Grouping (E.24) by
sample column gives \[
W_n-\mu_n=\frac1{\sqrt n}\sum_{i=1}^n V_i,\qquad
V_i=\eta\sum_{j\ge0}A^{-j-1}U(q,Z_{j,i}).
\tag{E.28}
\] The \(V_i\) are iid, centered, bounded by \(\eta B_U/(A-1)\), and
have variance \(s^2>0\). Berry--Esseen therefore gives \[
\sup_z\left|
\Pr(W_n\le z)-\Phi((z-\mu_n)/s)
\right|\le C/\sqrt n.
\tag{E.29}
\] For closed intervals the exact stopped-coordinate identities are \[
\{\tau_n\le t,J_n=+1\}
=\{Y_{n,t}>\sqrt n A^{-t}B_n^+\},
\] \[
\{\tau_n\le t,J_n=-1\}
=\{Y_{n,t}<-\sqrt n A^{-t}B_n^-\}.
\tag{E.30}
\] Before exit, the threshold cannot be crossed. After an exit at
\(j\le t\), the coordinate is frozen at a value whose magnitude is
larger than the boundary scaled by \(A^{-t}\).

For an excluded endpoint, replace the corresponding strict comparison in
(E.30) by a weak comparison. The smoothing argument below works for
either convention because the comparison Gaussian has a continuous
bounded density. No non-atomicity of the finite-batch updates is
assumed.

For any threshold, (E.27), Chebyshev's inequality, and (E.29) bound the
probability error by \[
C\{n^{-1/2}+h+\Delta_t/h^2\},\qquad
\Delta_t=(1+\log n)/n+A^{-2t}.
\] Choosing \(h=\Delta_t^{1/3}\) proves (E.16). At exit, the sign of
\(Y_{n,\tau_n}\) is \(J_n\), which proves (E.19).

\textbf{Step 5: remove the early-time loss in the bound.} The martingale
in (E.22) has a uniform sub-Gaussian proxy. Indeed, each increment has
conditional proxy at most \[
\eta^2 A^{-2j-2}B_U^2\sup_{\mathcal K}|H'|^2.
\] The sum is finite, so for every deterministic \(t\), \[
\mathbb E e^{u\mathcal M_{n,t}}\le e^{C u^2}.
\tag{E.31}
\] For sufficiently large \(y\), (E.22) and (E.31) imply \[
\Pr(|Y_{n,t}|>y)\le C e^{-c y^2}+\frac{C}{ny^2}.
\tag{E.32}
\] Uniform boundedness of \(\mu_n\) is used here.

If \(nA^{-2t}>D\log n\), both true exit probabilities in (E.30) and
their Gaussian approximations are at most \[
C e^{-cD\log n}+\frac{C}{nD\log n}.
\] Choose a fixed sufficiently large \(D\). On all remaining times,
\(A^{-2t}\le D\log n/n\), so (E.16) is bounded by \(Ce_n\). This proves
(E.18). \(\square\)

The logarithmic mechanism and the Gaussian seed are classical
unstable-exit phenomena {[}E-R2, E-R3{]}. The extension here keeps the
actual batch-gradient noise, a declared exit interval, discrete update
times, and the geometric quantities needed for the calibration analysis
below.

\subsection{References}\label{scalar_exit-references}

{[}E-R1{]} A. W. van der Vaart. \emph{Asymptotic Statistics}. Cambridge
University Press, 1998. Chapter 5, ``M- and Z-Estimators.''
\href{https://doi.org/10.1017/CBO9780511802256.006}{Publisher chapter}.

{[}E-R2{]} Y. Bakhtin. \emph{Exit asymptotics for small diffusion about
an unstable equilibrium}. arXiv:math/0701569, 2007.
\href{https://arxiv.org/abs/math/0701569}{Preprint}. Journal DOI:
\href{https://doi.org/10.1016/j.spa.2007.06.003}{10.1016/j.spa.2007.06.003}.

{[}E-R3{]} Y. Bakhtin and A. Gaál. \emph{Exit time asymptotics for
dynamical systems with fast random switching near an unstable
equilibrium}. arXiv:1901.05513, 2019.
\href{https://arxiv.org/abs/1901.05513}{Preprint}. Journal DOI:
\href{https://doi.org/10.1142/S0219493720500264}{10.1142/S0219493720500264}.

{[}E-R4{]} P. Massart. \emph{The Tight Constant in the
Dvoretzky--Kiefer--Wolfowitz Inequality}. The Annals of Probability
18(3):1269--1283, 1990.
\href{https://doi.org/10.1214/aop/1176990746}{DOI}.
